% !PRINTED 620
\weberpage{620}
%<<CONT
gehört, auch in $C$ enthalten sein. $C$ ist also ein Normaltheiler
von $P$, und folglich muss nach dem Theorem I. die Gruppe $P$
linear sein.

VII.~\gesperrt{Eine irreducible Gleichung von Primzahlgrad,
bei der alle Wurzeln rational durch zwei beliebige von ihnen
ausdrückbar sind, ist metacyklisch.}

Wir schliessen aus diesen Sätzen noch auf eine merkwürdige,
zuerst von \gesperrt{Kronecker} bemerkte Eigenschaft der metacyklischen
Gleichungen für den Fall eines reellen Rationalitätsbereichs\footnote{Ueber algebraisch auflösbare Gleichungen. Monatsbericht der Berliner
Akademie, 14. April 1856.}.
Wenn bei einer solchen Gleichung zwei Wurzeln reell sind, so
folgt aus VI, dass alle ihre Wurzeln reell sind. Eine reelle
Wurzel muss aber eine solche Gleichung, da sie ungeraden
Grades ist, immer haben. Also folgt

VIII.~\gesperrt{Eine irreducible metacyklische Gleichung von
ungeradem Primzahlgrad mit reellen Co\"efficienten hat entweder
lauter reelle Wurzeln oder nur eine.}

Sind alle Wurzeln reell, so ist die Discriminante als Product
von lauter Quadraten reeller Grössen $(x_h - x_k)^2$ positiv. Ist nur
eine Wurzel reell, so entspricht jedem complexen Factor der
Discriminante $(x_h - x_k)^2$ ein conjugirter, und deren Product ist
positiv. Nur wenn $x_h$ und $x_k$ conjugirt imaginär, also $x_h - x_k$
rein imaginär ist, so ist $(x_h - x_k)^2$ negativ und die Discriminante
erhält also für jedes Paar conjugirt imaginärer Wurzeln den
Factor $-1$. Das Vorzeichen der Discriminante ist also $(-1)^{\frac{n-1}{2}}$.
Daraus folgt:

IX.~\gesperrt{Wenn $n \equiv 1\ (\mathrm{mod}\ 4)$ ist, so ist die
Discriminante immer positiv, ist aber $n \equiv 3\ (\mathrm{mod}\ 4)$, so
entscheidet das Vorzeichen der Discriminante, welcher der beiden
Fälle des Theorems VIII. eintritt.}
